\numoptsfalse
\yearpaper{2022}{Version D-1 \textperiodcentered{} 60 Questions \textperiodcentered{} 60 Marks \textperiodcentered{} 80 Minutes}{}
\begin{multicols}{2}

\QQ{In case of Fraunhoffer diffraction at a single slit the diffraction pattern on the screen is correct for which of the following statements?}
{\opt{Central dark band having uniform brightness on either side.}{Central dark band having alternate dark and bright bands of decreasing intensity on either side.}{Central bright band having dark bands on either side.}{Central bright band having alternate dark and bright bands of decreasing intensity on either side.}}

\QQ{When a Compact Disc (CD) is illuminated by small source of white light coloured bands are observed. This is due to}
{\opt{Interference}{Scattering}{Reflection}{Diffraction}}

\QQ{Consider a glass slab which is silvered at one side and the other side is transparent. Given the refractive index of the glass slab to be 1.5. If a ray of light is incident at an angle of $45^\circ$ on the transparent side, the deviation of the ray of light from its initial path, when it comes out of the slab is}
{\opt{$120^\circ$}{$90^\circ$}{$45^\circ$}{$180^\circ$}}

\QQ{Focal length of a convex lens will be maximum for}
{\opt{Green light}{Blue light}{Red light}{Yellow light}}

\QQ{For light diverging from a finite point source}
{\opt{the wave front is parabolic.}{the wave front is cylindrical}{the intensity at the wave front does not depend on the distance.}{the intensity decreases in proportion to the distance squared.}}

\QQ{The fringe width for red colour as compared to that for violet colour is approximately}
{\opt{4 times}{3 times}{8 times}{Double}}

\QQ{In accordance with the Bohr's model, the quantum number that characterises the Earth's revolution around the Sun in an orbit of radius $1.5\times10^{11}\unit{m}$ with orbital speed $3\times10^4\unit{m\,s^{-1}}$ is [Given mass of Earth $= 6\times10^{24}\unit{kg}$]}
{\opt{$8.57\times10^{64}$}{$5.98\times10^{86}$}{$2.57\times10^{74}$}{$2.57\times10^{38}$}}

\QQ{An electron is revolving in its Bohr orbit having Bohr radius of $0.529\unit{\text{\AA}}$, then the radius of third orbit is}
{\opt{$4.761\unit{\text{\AA}}$}{$4234\unit{nm}$}{$5125\unit{nm}$}{$4496\unit{\text{\AA}}$}}

\QQ{Binding energy of a Nitrogen nucleus $\left[{}^{14}_{7}\mathrm N\right]$, given $m\left[{}^{14}_{7}\mathrm N\right] = 14.00307\unit{u}$}
{\opt{$206.5\unit{MeV}$}{$104.7\unit{MeV}$}{$78\unit{MeV}$}{$85\unit{MeV}$}}

\QQ{In a photo electric experiment, if both the intensity and frequency of the incident light are doubled, then the saturation photo electric current}
{\opt{is doubled}{remains constant}{becomes four times}{is halved}}

\QQ{The kinetic energy of the photoelectrons increases by $0.52\unit{eV}$ when the wavelength of incident light is changed from $500\unit{nm}$ to another wavelength which is approximately}
{\opt{$1250\unit{nm}$}{$700\unit{nm}$}{$1000\unit{nm}$}{$400\unit{nm}$}}

\QQ{The de-Broglie wavelength of a particle of kinetic energy `$K$' is $\lambda$; the wavelength of the particle, if its kinetic energy is $\dfrac K4$ is}
{\opt{$\dfrac\lambda2$}{$\lambda$}{$4\lambda$}{$2\lambda$}}

\QQ{The radius of hydrogen atom in the ground state is $0.53\unit{\text{\AA}}$. After collision with an electron, it is found to have a radius of $2.12\unit{\text{\AA}}$, the principal quantum number `$n$' of the final state of the atom is}
{\opt{$n=3$}{$n=1$}{$n=4$}{$n=2$}}

\QQ{Which logic gate is represented by the following combination of logic gates?
\fig{\begin{circuitikz}[scale=0.7,transform shape]
\draw (2.6,0.6) node[nor port](G){};
\draw (0,1.2) node[not port,anchor=in](N1){}; \draw (0,0) node[not port,anchor=in](N2){};
\draw (N1.in) -- ++(-0.4,0) node[circ]{} node[left]{A}; \draw (N2.in) -- ++(-0.4,0) node[circ]{} node[left]{B};
\draw (N1.out) -| ($(G.in 1)+(-0.3,0)$) -- (G.in 1); \draw (N2.out) -| ($(G.in 2)+(-0.3,0)$) -- (G.in 2);
\draw (G.out) -- ++(0.4,0) node[right]{Y};
\end{circuitikz}}}
{\opt{AND}{OR}{NOR}{NAND}}

\QQ{A metallic rod of mass per unit length $0.5\unit{kg\,m^{-1}}$ is lying horizontally on a smooth inclined plane which makes an angle of $30^\circ$ with the horizontal. A magnetic field of strength $0.25\unit{T}$ is acting on it in the vertical direction. When a current `$I$' is flowing through it, the rod is not allowed to slide down. The quantity of current required to keep the rod stationary is}
{\opt{$14.76\unit{A}$}{$7.14\unit{A}$}{$11.32\unit{A}$}{$5.98\unit{A}$}}

\QQ{A nuclear reactor delivers a power of $10^9\unit{W}$, the amount of fuel consumed by the reactor in one hour is}
{\opt{$0.72\unit{g}$}{$0.04\unit{g}$}{$0.96\unit{g}$}{$0.08\unit{g}$}}

\QQ{Which of the following radiations is deflected by electric field?}
{\opt{$\gamma$-rays}{X-rays}{$\alpha$-particles}{Neutrons}}

\QQ{The resistivity of a semiconductor at room temperature is in between}
{\opt{$10^6$ to $10^8\unit{\Omega\,cm}$}{$10^{-2}$ to $10^{-5}\unit{\Omega\,cm}$}{$10^{10}$ to $10^{12}\unit{\Omega\,cm}$}{$10^{-3}$ to $10^{6}\unit{\Omega\,cm}$}}

\QQ{The forbidden energy gap for `Ge' crystal at `0' K is}
{\opt{$2.57\unit{eV}$}{$0.071\unit{eV}$}{$6.57\unit{eV}$}{$0.71\unit{eV}$}}

\QQ{Two masses of $5\unit{kg}$ and $3\unit{kg}$ are suspended with the help of massless inextensible strings as shown in fig. When whole system is going upwards with acceleration $2\unit{m/s^2}$, the value of $T_1$ is (use $g = 9.8\unit{m/s^2}$)
\fig{\begin{tikzpicture}[scale=0.75]
\fill[pattern=north east lines] (-1,2.6) rectangle (1,2.75); \draw (-1,2.6) -- (1,2.6);
\draw (0,2.6) -- (0,1.7); \draw[->] (0,2.0) -- (0,2.3) node[right,lbl]{$T_1$};
\draw[pattern=north east lines] (-0.3,1.1) rectangle (0.3,1.7); \node[lbl,right] at (0.35,1.4) {5 kg};
\draw (0,1.1) -- (0,0.5); \draw[->] (0,0.7) -- (0,1.0) node[right,lbl]{$T_2$};
\draw[pattern=north east lines] (-0.3,0) rectangle (0.3,0.5); \node[lbl,right] at (0.35,0.25) {3 kg};
\end{tikzpicture}}}
{\opt{$23.6\unit{N}$}{$94.4\unit{N}$}{$59\unit{N}$}{$35.4\unit{N}$}}

\QQ{The Vernier scale of a travelling microscope has 50 divisions which coincides with 49 main scale divisions. If each main scale division is $0.5\unit{mm}$, then the least count of the microscope is}
{\opt{$0.01\unit{mm}$}{$0.01\unit{cm}$}{$0.5\unit{cm}$}{$0.5\unit{mm}$}}

\QQ{The displacement `$x$' (in metre) of a particle of mass `$m$' (in kg) moving in one dimension under the action of a force, is related to time `$t$' (in sec.) by, $t = \sqrt x+3$. The displacement of the particle when its velocity is zero, will be}
{\opt{$6\unit{m}$}{$4\unit{m}$}{$2\unit{m}$}{$0\unit{m}$}}

\QQ{Two objects are projected at an angle $\theta^\circ$ and $(90-\theta)^\circ$, to the horizontal with the same speed. The ratio of their maximum vertical heights is}
{\opt{$1:\tan\theta$}{$1:1$}{$\tan^2\theta:1$}{$\tan\theta:1$}}

\QQ{A car is moving in a circular horizontal track of radius $10\unit{m}$ with a constant speed of $10\unit{m\,s^{-1}}$. A bob is suspended from the roof of the car by a light wire of length $1.0\unit{m}$. The angle made by the wire with the vertical is (in radian)}
{\opt{0}{$\dfrac\pi6$}{$\dfrac\pi3$}{$\dfrac\pi4$}}

\QQ{A tiny spherical oil drop carrying a net charge $q$ is balanced in still air, with a vertical uniform electric field of strength $\dfrac{81}{7}\pi\times10^5\unit{V/m}$. When the field is switched off, the drop is observed to fall with terminal velocity $2\times10^{-3}\unit{m\,s^{-1}}$. Here $g = 9.8\unit{m/s^2}$, viscosity of air is $1.8\times10^{-5}\unit{N\,s/m^2}$ and the density of oil is $900\unit{kg\,m^{-3}}$. The magnitude of `$q$' is}
{\opt{$1.6\times10^{-19}\unit{C}$}{$0.8\times10^{-19}\unit{C}$}{$3.2\times10^{-19}\unit{C}$}{$8\times10^{-19}\unit{C}$}}

\QQ{``Heat cannot be itself flow from a body at lower temperature to a body at higher temperature''. This statement corresponds to}
{\opt{Conservation of mass}{Second law of Thermodynamics}{First law of thermodynamics}{Conservation of momentum}}

\QQ{A smooth chain of length $2\unit{m}$ is kept on a table such that its length of $60\unit{cm}$ hangs freely from the edge of the table. The total mass of the chain is $4\unit{kg}$. The work done in pulling the entire chain on the table is, (Take $g = 10\unit{m/s^2}$)}
{\opt{$3.6\unit{J}$}{$12.9\unit{J}$}{$2.0\unit{J}$}{$6.3\unit{J}$}}

\QQ{The angular speed of a motor wheel is increased from $1200\unit{rpm}$ to $3120\unit{rpm}$ in 16 seconds. The angular acceleration of the motor wheel is}
{\opt{$6\pi\unit{rad/s^2}$}{$2\pi\unit{rad/s^2}$}{$8\pi\unit{rad/s^2}$}{$4\pi\unit{rad/s^2}$}}

\QQ{The centre of mass of an extended body on the surface of the earth and its centre of gravity}
{\opt{can never be at the same point.}{are always at the same point for any size of the body.}{centre of mass coincides with the centre of gravity of a body if the size of the body is negligible as compared to the size (or radius) of the earth.}{are always at the same point only for spherical bodies.}}

\QQ{A metallic rod breaks when strain produced is 0.2\%. The Young's modulus of the material of the rod is $7\times10^9\unit{N/m^2}$. The area of cross section to support a load of $10^4\unit{N}$ is}
{\opt{$7.1\times10^{-4}\unit{m^2}$}{$7.1\times10^{-8}\unit{m^2}$}{$7.1\times10^{-2}\unit{m^2}$}{$7.1\times10^{-6}\unit{m^2}$}}

\QQ{A charged particle of mass `$m$' and charge `$q$' is released from rest in an uniform electric field $\vec E$. Neglecting the effect of gravity, the kinetic energy of the charged particle after `$t$' second is}
{\opt{$\dfrac{Eqm}{t}$}{$\dfrac{2E^2t^2}{mq}$}{$\dfrac{E^2q^2t^2}{2m}$}{$\dfrac{Eq^2m}{2t^2}$}}

\QQ{The electric field and the potential of an electric dipole vary with distance $r$ as}
{\opt{$\dfrac1{r^2}$ and $\dfrac1{r^3}$}{$\dfrac1r$ and $\dfrac1{r^2}$}{$\dfrac1{r^3}$ and $\dfrac1{r^2}$}{$\dfrac1{r^2}$ and $\dfrac1r$}}

\QQ{The displacement of a particle executing SHM is given by $x = 3\sin\left[2\pi t+\dfrac\pi4\right]$ where `$x$' is in metres and `$t$' is in seconds. The amplitude and maximum speed of the particle is}
{\opt{$3\unit{m}$, $6\pi\unit{m\,s^{-1}}$}{$3\unit{m}$, $2\pi\unit{m\,s^{-1}}$}{$3\unit{m}$, $8\pi\unit{m\,s^{-1}}$}{$3\unit{m}$, $4\pi\unit{m\,s^{-1}}$}}

\QQ{Electrical as well as gravitational affects can be thought to be caused by fields. Which of the following is true for an electrical or gravitational field?}
{\opt{Fields are useful for understanding forces acting through a distance.}{The field concept is often used to describe contact forces.}{There is no way to verify the existence of a force field since it is just a concept.}{Gravitational or Electric field does not exist in the space around an object.}}

\QQ{Four charges $+q$, $+2q$, $+q$ and $-2q$ are placed at the corners of a square ABCD respectively. The force on a unit positive charge kept at the centre `O' is}
{\opt{along the diagonal AC}{zero}{perpendicular to AD}{along the diagonal BD}}

\QQ{An electric dipole with dipole moment $4\times10^{-9}\unit{C\,m}$ is aligned at $30^\circ$ with the direction of a uniform electric field of magnitude $5\times10^4\unit{N\,C^{-1}}$, the magnitude of the torque acting on the dipole is}
{\opt{$10^{-5}\unit{N\,m}$}{$10^{-4}\unit{N\,m}$}{$10\times10^{-3}\unit{N\,m}$}{$\sqrt3\times10^{-4}\unit{N\,m}$}}

\QQ{Ten identical cells each of potential `$E$' and internal resistance `$r$', are connected in series to form a closed circuit. An ideal voltmeter connected across three cells, will read}
{\opt{$13E$}{$10E$}{$7E$}{$3E$}}

\QQ{In an atom electrons revolve around the nucleus along a path of radius $0.72\unit{\text{\AA}}$ making $9.4\times10^{18}$ revolutions per second. The equivalent current is -- [given $e = 1.6\times10^{-19}\unit{C}$]}
{\opt{$1.4\unit{A}$}{$1.2\unit{A}$}{$1.8\unit{A}$}{$1.5\unit{A}$}}

\QQ{When a metal conductor connected to left gap of a meter bridge is heated, the balancing point}
{\opt{remains unchanged}{shifts towards right}{shifts to the center}{shifts towards left}}

\QQ{Two tiny spheres carrying charges $1.8\unit{\mu C}$ and $2.8\unit{\mu C}$ are located at $40\unit{cm}$ apart. The potential at the mid-point of the line joining the two charges is}
{\opt{$4.3\times10^4\unit{V}$}{$3.8\times10^4\unit{V}$}{$3.6\times10^5\unit{V}$}{$2.1\times10^5\unit{V}$}}

\QQ{A parallel plate capacitor is charged by connecting a $2\unit{V}$ battery across it. It is then disconnected from the battery and a glass slab is introduced between plates. Which of the following pairs of quantities decrease?}
{\opt{Energy stored and capacitance.}{Charge and potential difference.}{Capacitance and charge.}{Potential difference and energy stored.}}

\QQ{A charged particle is moving in an electric field of $3\times10^{-10}\unit{V\,m^{-1}}$ with mobility $2.5\times10^6\unit{m^2/V/s}$, its drift velocity is}
{\opt{$2.5\times10^4\unit{m/s}$}{$7.5\times10^{-4}\unit{m/s}$}{$1.2\times10^{-4}\unit{m/s}$}{$8.33\times10^{-4}\unit{m/s}$}}

\QQ{Wire bound resistors are made by}
{\opt{winding the wires of an alloy of Ge, Au, Ga}{winding the wires of an alloy of Cu, Al, Ag}{winding the wires of an alloy of manganin, constantan, nichrome}{winding the wires of an alloy of Si, Tu, Fe}}

\QQ{A galvanometer of resistance $50\unit{\Omega}$ is connected to a battery of $3\unit{V}$ along with a resistance $2950\unit{\Omega}$ in series. A full scale deflection of 30 divisions is obtained in the galvanometer. In order to reduce this deflection to 20 divisions, the resistance in series should be}
{\opt{$5050\unit{\Omega}$}{$6050\unit{\Omega}$}{$4450\unit{\Omega}$}{$5550\unit{\Omega}$}}

\QQ{A circular coil of wire of radius `$r$' has `$n$' turns and carries a current `$I$'. The magnetic induction `$B$' at a point on the axis of the coil at a distance $\sqrt3\,r$ from its centre is}
{\opt{$\dfrac{\mu_0nI}{16r}$}{$\dfrac{\mu_0nI}{32r}$}{$\dfrac{\mu_0nI}{4r}$}{$\dfrac{\mu_0nI}{8r}$}}

\QQ{If voltage across a bulb rated $220\unit{V}$, $100\unit{W}$ drops by 2.5\% of its rated value, the percentage of the rated value by which the power would decrease is}
{\opt{5\%}{20\%}{10\%}{2.5\%}}

\QQ{A wire of a certain material is stretched slowly by 10\%. Its new resistance and specific resistance becomes respectively}
{\opt{1.21 times, same}{1.1 times, 1.1 times}{both remains the same}{1.2 times, 1.1 times}}

\QQ{A proton moves with a velocity of $5\times10^6\,\hat j\unit{m\,s^{-1}}$ through the uniform electric field, $\vec E = 4\times10^6\,[2\hat i+0.2\hat j+0.1\hat k]\unit{V\,m^{-1}}$ and the uniform magnetic field $\vec B = 0.2\,[\hat i+0.2\hat j+\hat k]\unit{T}$. The approximate net force acting on the proton is}
{\opt{$2.2\times10^{-13}\unit{N}$}{$5\times10^{-13}\unit{N}$}{$20\times10^{-13}\unit{N}$}{$25\times10^{-13}\unit{N}$}}

\QQ{A solenoid of length $50\unit{cm}$ having 100 turns carries a current of $2.5\unit{A}$. The magnetic field at one end of the solenoid is}
{\opt{$1.57\times10^{-4}\unit{T}$}{$3.14\times10^{-4}\unit{T}$}{$9.42\times10^{-4}\unit{T}$}{$6.28\times10^{-4}\unit{T}$}}

\QQ{An alternating current is given by $i = i_1\sin\omega t+i_2\cos\omega t$. The r.m.s.\ current is given by}
{\opt{$\sqrt{\dfrac{i_1^2+i_2^2}{2}}$}{$\dfrac{i_1+i_2}{\sqrt2}$}{$\sqrt{\dfrac{i_1^2+i_2^2}{\sqrt2}}$}{$\dfrac{i_1-i_2}{\sqrt2}$}}

\QQ{Which of the following statements proves that Earth has a magnetic field?}
{\opt{Earth is surrounded by ionosphere.}{The intensity of cosmic rays stream of charged particles is more at the poles than at the equator.}{A large quantity of iron-ore is found in the Earth.}{Earth is a planet rotating about the North South axis.}}

\QQ{A long solenoid has 500 turns, when a current of $2\unit{A}$ is passed through it, the resulting magnetic flux linked with each turn of the solenoid is $4\times10^{-3}\unit{Wb}$, then self induction of the solenoid is}
{\opt{2.0 henry}{4.0 henry}{1.0 henry}{2.5 henry}}

\QQ{A fully charged capacitor `$C$' with initial charge `$q_0$' is connected to a coil of self inductance `$L$' at $t = 0$. The time at which the energy is stored equally between the electric and the magnetic field is}
{\opt{$\pi\sqrt{LC}$}{$2\pi\sqrt{LC}$}{$\dfrac\pi4\sqrt{LC}$}{$\sqrt{LC}$}}

\QQ{A magnetic field of flux density $1.0\unit{Wb\,m^{-2}}$ acts normal to a 80 turn coil of $0.01\unit{m^2}$ area. If this coil is removed from the field in 0.2 second, the emf induced in it is}
{\opt{$0.8\unit{V}$}{$4\unit{V}$}{$5\unit{V}$}{$8\unit{V}$}}

\QQ{A ray of light passes through an equilateral glass prism in such a manner that the angle of incidence is equal to the angle of emergence and each of these angles is equal to $\frac34$ of the angle of the prism. The angle of deviation is}
{\opt{$20^\circ$}{$45^\circ$}{$30^\circ$}{$39^\circ$}}

\QQ{A convex lens of focal length `$f$' is placed somewhere in between an object and a screen. The distance between the object and the screen is `$x$'. If the numerical value of the magnification produced by the lens is `$m$', then the focal length of the lens is}
{\opt{$\dfrac{(m+1)^2x}{m}$}{$\dfrac{mx}{(m+1)^2}$}{$\dfrac{(m-1)^2x}{m}$}{$\dfrac{mx}{(m-1)^2}$}}

\QQ{A series resonant ac circuit contains a capacitance $10^{-6}\unit{F}$ and an inductor of $10^{-4}\unit{H}$. The frequency of electrical oscillations will be}
{\opt{$\dfrac{10^5}{2\pi}\unit{Hz}$}{$10^5\unit{Hz}$}{$\dfrac{10}{2\pi}\unit{Hz}$}{$10\unit{Hz}$}}

\QQ{In a series LCR circuit $R = 300\unit{\Omega}$, $L = 0.9\unit{H}$, $C = 2.0\unit{\mu F}$ and $\omega = 1000\unit{rad/sec}$, then impedance of the circuit is}
{\opt{$500\unit{\Omega}$}{$1300\unit{\Omega}$}{$400\unit{\Omega}$}{$900\unit{\Omega}$}}

\QQ{Which of the following radiations of electromagnetic waves has the highest wavelength?}
{\opt{IR-rays}{X-rays}{Microwaves}{UV-rays}}

\QQ{The power of a equi-concave lens is $-4.5\unit{D}$ and is made of a material of R.I.\ 1.6, the radii of curvature of the lens is}
{\opt{$-2.66\unit{cm}$}{$-26.6\unit{cm}$}{$115.44\unit{cm}$}{$+36.6\unit{cm}$}}

\end{multicols}
